\documentclass[11pt]{article}
\usepackage[margin=1in]{geometry}
\usepackage[T1]{fontenc}
\usepackage{lmodern}
\usepackage{amsmath,amssymb,amsthm,mathtools}
\usepackage{graphicx}
\usepackage{xcolor}
\usepackage[hidelinks]{hyperref}

\newtheorem{theorem}{Theorem}
\newtheorem{lemma}{Lemma}
\newtheorem{remark}{Remark}
\newcommand{\R}{\mathbb R}
\newcommand{\C}{\mathbb C}
\newcommand{\T}{\mathbb T}
\newcommand{\N}{\mathbb N}
\newcommand{\F}{\mathcal F}
\newcommand{\SN}{\mathring S}

\title{Spherical Hermite sums on modulation spaces\\
       \large A complete classification in dimensions at least three}
\author{Codex agent lane 6}
\date{August 11, 2026}

\begin{document}
\maketitle

\begin{center}
\fbox{\parbox{0.91\linewidth}{
\textbf{Status: candidate partial result, likely valid.}
The source question is completely answered for every $d\ge3$, but the
intermediate range in $d=2$ remains open.  Human review is recommended.}}
\end{center}

\section*{Source question}

Jaming and Speckbacher ask in Remark 3.3 of
\emph{Convergence of Hermite expansions in modulation spaces},
arXiv:2601.04893 \cite{JS}, for which $p$ the spherical Hermite sums
\[
  \SN_N f=\sum_{\substack{n\in\N_0^d\\ |n|\le N}}
       \langle f,h_n\rangle h_n
\]
converge in $M^p(\R^d)$.  Here $h_n$ are the tensor Hermite functions and

\begin{center}
\includegraphics[width=0.96\linewidth]{figures/open_problem_crop.png}
\end{center}

The image is a full-width crop of source page 10.  The question follows the
authors' observation that the known implication from torus Fourier multipliers
to Hermite multipliers is only sufficient; they specifically point to
Fefferman's ball multiplier as the natural obstruction.

\section*{Result and scope}

\begin{theorem}[High-dimensional classification]
Let $d\ge3$ and $1\le p\le\infty$.  The following are equivalent:
\begin{enumerate}
  \item for every $f\in M^p(\R^d)$, one has
        $\SN_Nf\to f$ in $M^p(\R^d)$;
  \item $p=2$.
\end{enumerate}
In particular, for every $p\ne2$,
\[
  \sup_{N\in\N}\|\SN_N\|_{M^p\to M^p}=\infty.
\]
For $d=2$, convergence holds at $p=2$ and fails at $p=1,\infty$; the
present proof does not decide $1<p<\infty$, $p\ne2$.
\end{theorem}

\section*{Proof intuition}

The Bargmann transform sends Hermite functions to normalized monomials.  A
finite collection of monomials whose multi-indices lie near a very large base
index $K$ has, after division by the norm of the base monomial, the same
asymptotic norm as the corresponding trigonometric polynomial on a torus.  This
is an elementary concentration fact: the radial variables under the mass of a
large monomial are Gamma distributed and concentrate at the base index.

We place the neighboring indices in the integer hyperplane $v^\perp$.  Pell's
equation produces arbitrarily remote base indices $K$ and integer radii $N$
for which the spherical condition is \emph{exactly}
$|q|\le J$ on that hyperplane.  Thus a hypothetical uniform bound for
spherical Hermite sums would imply uniform bounds for ellipsoidal Fourier sums
on a $(d-1)$-torus.  When $d\ge3$, de Leeuw transference and Fefferman's
theorem make this impossible unless $p=2$.

\section*{Bargmann-Fock setup}

Under the standard Bargmann transform, $M^p(\R^d)$ is isomorphic (isometric
up to the normalization convention) to the Fock space
\[
 \F^p(\C^d)=\left\{F\text{ entire}:
 \int_{\C^d}|F(z)|^p e^{-\pi p|z|^2/2}\,dz<\infty\right\},
 \qquad 1\le p<\infty,
\]
and $h_n$ is sent to
\[
 e_n(z)=\frac{\pi^{|n|/2}}{\sqrt{n!}}z^n,
 \qquad n\in\N_0^d.
\]
Consequently $\SN_N$ becomes the monomial multiplier with symbol
$1_{\{|n|\le N\}}$.

\begin{lemma}[High-monomial torus limit]\label{lem:focklimit}
Fix $1\le p<\infty$, a vector $v\in\N^d$ all of whose coordinates are
positive, a finite set
\[
 E\subset H_v:=\{q\in\mathbb Z^d:v\cdot q=0\},
\]
and scalars $c_q$.  For $K=Lv$, with $L$ sufficiently large that
$K+q\in\N_0^d$ for every $q\in E$, put
\[
 F_L=\sum_{q\in E}c_qe_{K+q}.
\]
Then
\[
 \lim_{L\to\infty}\frac{\|F_L\|_{\F^p}}{\|e_K\|_{\F^p}}
 =\left\|\sum_{q\in E}c_q e^{iq\cdot\theta}\right\|_{L^p(\T^d)},
\]
where Haar measure on the torus is normalized.
\end{lemma}

\begin{proof}
Write $z_j=r_je^{i\theta_j}$ and $t_j=\pi r_j^2$.  Normalize the measure

\[
 |e_K(z)|^p e^{-\pi p|z|^2/2}\,dz
\]

to have mass one.  Its angular part is normalized Haar measure, and its radial
variables are independent with
\[
 t_j\sim\operatorname{Gamma}\left(\frac{pK_j}{2}+1,
                                  \text{ rate }\frac p2\right).
\]
Furthermore,
\[
 \frac{e_{K+q}(z)}{e_K(z)}
 =A_{K,q}(t)e^{iq\cdot\theta},\qquad
 A_{K,q}(t)=\prod_{j=1}^d t_j^{q_j/2}
            \sqrt{\frac{K_j!}{(K_j+q_j)!}}.
\]
For fixed $q$, $t_j/K_j\to1$ in probability.  The factorial ratio
asymptotic therefore gives $A_{K,q}\to1$ in probability.  This convergence
also holds in every fixed finite $L^s$.  Indeed, for all sufficiently large
$K_j$, the Gamma moment formula gives
\begin{align*}
 \mathbb E[A_{K,q}^{s}]
 =\prod_{j=1}^d&\left(\frac{K_j!}{(K_j+q_j)!}\right)^{s/2}
 \left(\frac p2\right)^{-sq_j/2}\\
 &\times
 \frac{\Gamma(pK_j/2+1+sq_j/2)}{\Gamma(pK_j/2+1)}
 \longrightarrow1.
\end{align*}
Applying this with a moment strictly larger than $p$ supplies uniform
integrability, so $A_{K,q}\to1$ in $L^p$.  Since $E$ is finite,
Minkowski's inequality yields
\[
 \left\|\sum_{q\in E}c_q(A_{K,q}-1)e^{iq\cdot\theta}\right\|_{L^p}
 \le\sum_{q\in E}|c_q|\,\|A_{K,q}-1\|_{L^p}\longrightarrow0.
\]
This is exactly the asserted normalized norm limit.
\end{proof}

\section*{Exact integer-radius placement}

\begin{lemma}[Pell slice]\label{lem:pell}
For every $d\ge3$, there is a primitive vector $v\in\N^d$, all of whose
coordinates are positive, for which $D=|v|^2$ is not a square.  For each
integer $J\ge1$, there are sequences $K_k\in\N^d$ and $N_k\in\N$, with
$\min_j(K_k)_j\to\infty$, such that for every $q\in H_v$ and all large
$k$ for which $K_k+q\ge0$,
\[
 |K_k+q|\le N_k \quad\Longleftrightarrow\quad |q|\le J.
\]
\end{lemma}

\begin{proof}
If $d$ is not a square, take $v=(1,\ldots,1)$, so $D=d$.  If $d$ is a
square, take $v=(1,\ldots,1,2)$, with $d-1$ entries equal to one; then
$D=d+3$.  For square $d\ge4$, $d+3$ cannot be a square, because two
positive squares differ by three only in the pair $1,4$.  In either case
$v$ is primitive and $D$ is nonsquare.

Pell's equation
\[
 a^2-Db^2=1
\]
has positive solutions $(a_k,b_k)$ with $b_k\to\infty$.  Set
\[
 K_k=Jb_kv,\qquad N_k=Ja_k.
\]
For $q\in H_v$, orthogonality and Pell's identity give
\[
 |K_k+q|^2=DJ^2b_k^2+|q|^2=N_k^2-J^2+|q|^2.
\]
The desired equivalence follows, with no rounding or limiting error.
\end{proof}

\section*{Proof of the theorem}

Let first $1<p<\infty$, and suppose that $\SN_Nf\to f$ in $M^p$ for
every $f$.  Banach--Steinhaus yields
\[
 \sup_N\|\SN_N\|_{M^p\to M^p}=C<\infty.
\]
Transfer this bound to $\F^p(\C^d)$.  Choose $v$ from Lemma
\ref{lem:pell}.  Since $v$ is primitive, the lattice $H_v$ is saturated,
has rank $d-1$, and is the character lattice of a $(d-1)$-torus.

Fix $J\ge1$ and a finite trigonometric polynomial
\[
 P(\theta)=\sum_{q\in E}c_qe^{iq\cdot\theta},\qquad E\subset H_v.
\]
Embed it as $F_k=\sum_{q\in E}c_qe_{K_k+q}$, where $K_k,N_k$ come from
Lemma \ref{lem:pell}.  The exact Pell equivalence gives
\[
 \SN_{N_k}F_k=\sum_{\substack{q\in E\\|q|\le J}}c_qe_{K_k+q}.
\]
Apply the uniform Fock bound, divide by $\|e_{K_k}\|_{\F^p}$, and use Lemma
\ref{lem:focklimit} on both sides.  We obtain
\[
 \left\|\sum_{\substack{q\in E\\|q|\le J}}c_qe^{iq\cdot\theta}\right\|_p
 \le C\left\|\sum_{q\in E}c_qe^{iq\cdot\theta}\right\|_p. \tag{1}
\]

Choose a lattice isomorphism $B:\mathbb Z^{d-1}\to H_v$.  In these
coordinates, (1) says that the Fourier multipliers
\[
 m_J(\ell)=1_{\{Q(\ell)\le J^2\}},\qquad
 Q(\ell)=|B\ell|^2,
\]
are uniformly bounded on $L^p(\T^{d-1})$.  The dilation part of de Leeuw's
multiplier theorem \cite{deLeeuw} then makes the ellipsoid indicator
$1_{\{Q(\xi)\le1\}}$ an $L^p(\R^{d-1})$ Fourier multiplier.  The ellipsoid
boundary has measure zero, so the usual a.e.-continuity formulation applies.
After an invertible linear change of variables, this is the Euclidean ball
multiplier.  Since $d-1\ge2$, Fefferman's theorem \cite{Fefferman} forces
$p=2$, a contradiction whenever $p\ne2$.

At $p=1$, take $f=g\otimes h_0^{\otimes(d-1)}$.  Tensor factorization of
the modulation norm and the Hermite coefficients gives
\[
 \SN_N f=(S_Ng)\otimes h_0^{\otimes(d-1)},
\]
where $S_N$ is the one-dimensional Hermite partial sum.  The source proves
that $\|S_N\|_{M^1(\R)\to M^1(\R)}$ grows at least logarithmically
\cite[Theorem 1.1]{JS}; hence the spherical sums are not uniformly bounded on
$M^1(\R^d)$.  Under the standard $M^1$--$M^\infty$ duality, these
finite-rank real diagonal projections are adjoints of themselves, so their
$M^\infty$ operator norms equal their $M^1$ norms.  This rules out
$p=\infty$.

Finally, $M^2(\R^d)=L^2(\R^d)$ with equivalent (indeed, conventionally
isometric) norm by Moyal's identity.  The operators $\SN_N$ are increasing
orthogonal projections associated with the tensor Hermite orthonormal basis,
so $\SN_Nf\to f$ for every $f\in M^2$.  This proves the theorem.

\section*{Dimension two, verification, and limitations}

The proof above exposes the precise dimension threshold: $H_v$ has rank
$d-1$.  When $d=2$, it is one-dimensional, and its ball multipliers are
interval partial sums, uniformly bounded for $1<p<\infty$.  Thus this method
does not settle the intermediate two-dimensional range.  The tensor endpoint
argument and the $p=2$ argument remain valid, giving the scoped claims stated
in the theorem.

The included script checks Pell identities for representative dimensions and
radii and checks the convergence of the exact Gamma moments for positive and
negative shifts.  These computations are sanity checks, not parts of the
proof.

A bounded novelty search on August 11, 2026 covered the local run indexes, the
exact source title and question phrases, Hermite/Fock multiplier converses,
monomial ball sums in Fock spaces, and Pell-based multiplier transference.  It
found the source and general multiplier literature, but no later answer to
Remark 3.3 and no matching argument.  This is not an exhaustive novelty claim.

\medskip
\noindent\textbf{Human-review recommendation.}
Check Lemma \ref{lem:focklimit}'s normalization and the discontinuous-symbol
form of de Leeuw dilation.  The Pell placement is an exact integer identity.
The packet must remain classified as partial until the intermediate (d=2)
range is settled.

\begin{thebibliography}{9}

\bibitem{JS}
P. Jaming and M. Speckbacher,
\emph{Convergence of Hermite expansions in modulation spaces},
arXiv:2601.04893, 2026.

\bibitem{deLeeuw}
K. de Leeuw,
\emph{On $L_p$ multipliers},
Ann. of Math. (2) \textbf{81} (1965), 364--379.

\bibitem{Fefferman}
C. Fefferman,
\emph{The multiplier problem for the ball},
Ann. of Math. (2) \textbf{94} (1971), 330--336.

\end{thebibliography}

\end{document}
